% =========================================================================
% APPENDIX B — STUDY PROTOCOLS
% =========================================================================

\chapter{Study Protocols}
\label{app:protocols}

This appendix documents the instruments of the empirical evaluation. The study is delivered as a self-guided online questionnaire, so the protocol is the questionnaire itself: its structure, its four variants, its comprehension items, and the consent and information text that precede it. Where the exact wording is still being finalised as part of the work in progress, this is marked explicitly; the structure around it is fixed.

% -------------------------------------------------------------------------
\section{Questionnaire Variants}
\label{app:b-variants}

The instrument exists in four variants, crossing two languages with two conditions:

\begin{itemize}
    \item \textbf{Language:} Italian and English, so that participants respond in their first language. Open responses are analysed in the language in which they are written (Appendix~\ref{app:limits-evaluate}).
    \item \textbf{Condition:} a \emph{control} version showing only the conventional static representation (C1), and an \emph{experimental} version showing the dynamic and experiential representation (C2 and C3).
\end{itemize}

\noindent The four variants (IT-control, IT-experimental, EN-control, EN-experimental) are identical in structure, background questions, and comprehension items. They differ only in language and in which representational sequences the participant is shown. This keeps the two groups comparable: the information conveyed is held constant, and only the representational form varies (Appendix~\ref{app:limits-evaluate}).

% -------------------------------------------------------------------------
\section{Structure of the Questionnaire}
\label{app:b-structure}

Each participant follows the same journey, in four stages.

\paragraph{1. Consent and information.}
The questionnaire opens with the plain-language information sheet and the consent statement (\S\ref{app:b-consent}). Participation proceeds only after consent is given.

\paragraph{2. Background questions.}
A short set of questions characterises the participant and allows prior knowledge to be accounted for in analysis:
\begin{itemize}
    \item age band;
    \item occupation or field of work or study;
    \item prior familiarity with the AMOC or with ocean-climate systems (a short self-rated scale, with an option to describe any prior knowledge).
\end{itemize}
These items collect no directly identifying information; responses are anonymised (Appendix~\ref{app:ethics-conduct}).

\paragraph{3. Guided exposure with interleaved comprehension.}
The participant is taken through the representational sequences of their assigned condition. Comprehension questions are placed \emph{between} sequences rather than only at the end, so that understanding is captured as it forms. The comprehension items are described in \S\ref{app:b-comprehension}.

\paragraph{4. Unassessed closing sequence.}
After the assessed part is complete, every participant --- including the control group --- is shown the fuller experiential sequence once, for interest. Responses to it are not assessed and do not enter the between-group comparison; it serves as a debrief and a return to the participant, and is placed outside the assessed journey by design (Appendix~\ref{app:limits-evaluate}).

% -------------------------------------------------------------------------
\section{Comprehension Items}
\label{app:b-comprehension}

The comprehension items are the study's primary evidence. They target the three properties the representation is designed to convey --- the system's dynamics, its non-linearity, and its uncertainty --- so that they measure understanding of what the artefact actually encodes rather than general climate knowledge.

The items lean on \emph{open-ended} questions, which reveal a participant's mental model rather than their ability to recognise a supplied answer, with \emph{closed} items possibly added in support to give a signal directly comparable between the two groups. The open responses are the primary material for the reflexive thematic analysis; any closed items are reported descriptively (Appendix~\ref{app:limits-evaluate}). The exact wording, number, and open/closed balance are finalised close to data collection, drafted and reviewed with the supervisor, and piloted before the study proper.

\vspace{0.5em}
\noindent\textit{[TO BE FINALISED: the frozen set of comprehension items, in both languages, with their placement between sequences. Draft items target: (i) how the system changes over time; (ii) whether its response is proportional or can shift abruptly; (iii) how confident one can be in what is shown.]}

% -------------------------------------------------------------------------
\section{Information Sheet and Consent}
\label{app:b-consent}

\textit{[TO BE FINALISED: the full participant information sheet and consent statement, in both languages. The information sheet states, in plain language: what the study involves and how long it takes; that participation is voluntary and may be stopped at any point without consequence; what data are collected (background questions and questionnaire responses) and that no directly identifying information is gathered; how responses are anonymised, stored securely under the GDPR, retained only as long as needed, and used only for this research; and contact details for the researcher and supervisor. Consent is given explicitly before the questionnaire begins.]}

The ethical commitments these documents implement are set out in Appendix~\ref{app:ethics-conduct}.